% =====================================================================
%  MCMP117 -- Fundamentals of Operating Systems
%  Week 8 Course Notes: Midterm Exam Sample Paper
%  Compile with XeLaTeX (twice):  xelatex MCMP117_Week8_Notes.tex
%  Needs ./figures/ and the ../code/ folder.
% =====================================================================
\documentclass[11pt,a4paper]{article}
\newcommand{\weeknum}{8}
% ---------- Page & fonts ---------------------------------------------
\usepackage[a4paper,margin=2.3cm,headheight=15pt]{geometry}
\usepackage{fontspec}
\IfFontExistsTF{Calibri}{\setmainfont{Calibri}}{\setmainfont{Carlito}}
\IfFontExistsTF{Cambria}{\newfontfamily\headfont{Cambria}}{\newfontfamily\headfont{Caladea}}
\IfFontExistsTF{Consolas}{\setmonofont[Scale=0.95]{Consolas}}{\setmonofont[Scale=0.88]{DejaVu Sans Mono}}
\usepackage{amsmath,amssymb}
\usepackage{microtype}
\linespread{1.12}
\setlength{\parskip}{0.55em}
\setlength{\parindent}{0pt}

% ---------- Colours (course design system) ----------------------------
\usepackage{xcolor}
\definecolor{dark}{HTML}{212121}
\definecolor{charcoal}{HTML}{36454F}
\definecolor{slate}{HTML}{50808E}
\definecolor{mint}{HTML}{02C39A}
\definecolor{mintdeep}{HTML}{028090}
\definecolor{light}{HTML}{F2F2F2}
\definecolor{muted}{HTML}{6B7679}

% ---------- Packages ---------------------------------------------------
\usepackage{graphicx}
\usepackage{float}
\usepackage{needspace}
\graphicspath{{figures/}}
\usepackage[font=small,labelfont={bf,color=mintdeep},justification=centering]{caption}
\usepackage{booktabs,longtable,array,tabularx}
\usepackage{enumitem}
\setlist{itemsep=0.25em,topsep=0.3em}
\usepackage{tikz}
\usetikzlibrary{arrows.meta,positioning,calc,shapes.geometric}
\usepackage{pgfplots}
\pgfplotsset{compat=1.17}
\usepackage[most]{tcolorbox}
\usepackage{titlesec}
\usepackage{fancyhdr}
\usepackage{xurl}
\usepackage[colorlinks=true,linkcolor=mintdeep,urlcolor=mintdeep,citecolor=mintdeep]{hyperref}

% ---------- Headings ----------------------------------------------------
\titleformat{\section}{\headfont\Large\bfseries\color{charcoal}}{\color{mint}\thesection}{0.7em}{}
\titleformat{\subsection}{\headfont\large\bfseries\color{charcoal}}{\color{mintdeep}\thesubsection}{0.6em}{}
\titleformat{\subsubsection}{\headfont\normalsize\bfseries\color{mintdeep}}{}{0em}{}
\titlespacing*{\section}{0pt}{1.6em}{0.6em}
\titlespacing*{\subsection}{0pt}{1.2em}{0.4em}

% ---------- Header / footer -----------------------------------------------
\pagestyle{fancy}
\fancyhf{}
\fancyhead[L]{\small\color{muted}MCMP117 --- Fundamentals of Operating Systems}
\fancyhead[R]{\small\color{muted}Week \weeknum{} Notes}
\fancyfoot[C]{\small\color{muted}\thepage}
\renewcommand{\headrulewidth}{0pt}

% ---------- Boxes ---------------------------------------------------------
\newtcolorbox{keyidea}[1][Key idea]{enhanced,breakable,colback=charcoal,colframe=charcoal,
  coltext=white,fonttitle=\headfont\bfseries\color{mint},title={#1},
  attach title to upper={\par\smallskip},boxrule=0pt,arc=3pt,left=10pt,right=10pt,top=8pt,bottom=8pt}
\newtcolorbox{example}[1][Worked example]{enhanced,colback=light,colframe=light,
  fonttitle=\headfont\bfseries\color{mintdeep},title={#1},attach title to upper={\par\smallskip},
  boxrule=0pt,arc=3pt,left=10pt,right=10pt,top=8pt,bottom=8pt}
\newtcolorbox{tryit}[1][Try it yourself]{enhanced,breakable,colback=white,colframe=mint,
  fonttitle=\headfont\bfseries\color{mintdeep},title={#1},attach title to upper={\par\smallskip},
  boxrule=0.8pt,arc=3pt,left=10pt,right=10pt,top=8pt,bottom=8pt}
\newtcolorbox{histnote}[1][Historical note]{enhanced,breakable,colback=white,colframe=slate!40,
  fonttitle=\headfont\bfseries\color{slate},title={#1},attach title to upper={\par\smallskip},
  boxrule=0.5pt,arc=3pt,left=10pt,right=10pt,top=8pt,bottom=8pt,fontupper=\small}

% Photo helper: fits image inside a box, never distorted (keepaspectratio)
\newcommand{\photo}[3]{\includegraphics[width=#1,height=#2,keepaspectratio]{#3}}

\newcommand{\term}[1]{\textbf{\color{mintdeep}#1}}
\newcommand{\cmd}[1]{\texttt{#1}}


% ---------- Code listings ----------------------------------------------------
\usepackage{listings}
\lstdefinestyle{course}{
  basicstyle=\ttfamily\small, columns=fullflexible, keepspaces=true,
  backgroundcolor=\color{dark}, rulecolor=\color{dark}, frame=single, framesep=6pt,
  xleftmargin=8pt, xrightmargin=8pt, aboveskip=0.8em, belowskip=0.6em,
  basicstyle=\ttfamily\small\color{light},
  keywordstyle=\color{mint}\bfseries, commentstyle=\color{white!60}\itshape,
  stringstyle=\color{orange!70!yellow}, numberstyle=\tiny\color{muted},
  numbers=left, numbersep=10pt, showstringspaces=false, breaklines=true,
  tabsize=4, upquote=true}
\lstdefinestyle{term}{style=course, numbers=none, keywordstyle=\color{light},
  commentstyle=\color{white!60}, stringstyle=\color{light}}
\lstset{style=course,language=C}
\usepackage{textcomp}

\begin{document}
% ===================== MCMP117 Midterm Exam -- SAMPLE PAPER (student version) =====================
% Built with the course preamble:
%   python3 ../../lib/assemble_notes.py en 8 midterm_sample_body.tex MCMP117_Midterm_Sample.tex "Midterm Exam Sample Paper"
% Every answer is verified in ../code/sample_exam.py (output ../code/out_sample.txt); see MCMP117_Midterm_AnswerKey.
\fancyhead[R]{\small\color{muted}Midterm Exam --- Sample Paper}
\setlength{\parskip}{0.35em}
\newcommand{\pts}[1]{\textcolor{mintdeep}{\textbf{[#1~pt]}}\enspace}
\newcommand{\ans}[1]{\par\smallskip\noindent\begin{tikzpicture}\draw[muted!60,rounded corners=2pt] (0,0) rectangle (\linewidth,#1);\end{tikzpicture}\par\medskip}
% empty time axis for a Gantt chart, 0..#1 ms
\newcommand{\ganttgrid}[2]{\par\noindent\begin{tikzpicture}[x=0.82cm]
\node[anchor=east,font=\small\bfseries] at (-0.2,0.31) {#2};
\draw[muted!70] (0,0) rectangle (#1,0.62);
\foreach \t in {0,...,#1} {\draw[muted!70] (\t,0) -- (\t,-0.1) node[below,font=\scriptsize,text=muted,inner sep=1pt] {\t};}
\end{tikzpicture}\par\smallskip}
% empty frame table: 12 references, 3 frames
\newcommand{\frametable}[1]{\begin{tabular}{@{}l|*{12}{>{\centering\arraybackslash}p{0.62cm}|}@{}}
\multicolumn{13}{@{}l}{\textbf{#1}}\\ \hline
ref & 1 & 0 & 4 & 2 & 1 & 0 & 1 & 2 & 3 & 1 & 2 & 3 \\ \hline
fr.\ 0 &&&&&&&&&&&& \\ \hline
fr.\ 1 &&&&&&&&&&&& \\ \hline
fr.\ 2 &&&&&&&&&&&& \\ \hline
fault &&&&&&&&&&&& \\ \hline
\end{tabular}}
% empty per-process table for C1
\newcommand{\proctable}{\begin{tabular}{|l|c|c|c|c|}
\hline
\textbf{Process} & \textbf{Finish} & \textbf{Turnaround} & \textbf{Waiting} & \textbf{Response} \\ \hline
A &&&& \\ \hline B &&&& \\ \hline C &&&& \\ \hline D &&&& \\ \hline
\textbf{Average} & --- &&& \\ \hline
\end{tabular}}

\thispagestyle{empty}
\begin{tcolorbox}[enhanced,colback=dark,colframe=dark,arc=0pt,boxrule=0pt,
  left=18pt,right=18pt,top=14pt,bottom=14pt,coltext=white]
{\small\color{mint}\textsc{MCMP117 \,·\, Fundamentals of Operating Systems}}\par\smallskip
{\headfont\fontsize{22}{26}\selectfont\bfseries Midterm Exam --- Sample Paper}\par\smallskip
{\small Istanbul Okan University \,·\, Vocational School \,·\, Electronical Communication Technology Programme}\par
{\small Instructor: Asst.\ Prof.\ Dr.\ Hasan Oğuz \quad Date: [to be announced] \quad Room: [to be announced] \quad Duration: [90] minutes}
\end{tcolorbox}

\begin{center}
\renewcommand{\arraystretch}{1.6}
\begin{tabularx}{\linewidth}{@{}lX lX@{}}
\textbf{Name, surname:} & \hrulefill & \textbf{Student ID:} & \hrulefill \\
\textbf{Signature:} & \hrulefill & \textbf{Seat / room:} & \hrulefill \\
\end{tabularx}

\vspace{0.4em}
\renewcommand{\arraystretch}{1.35}
\begin{tabular}{|l|c|c|c|c|c|c|c|c|}
\hline
\textbf{Part} & A & B & C1 & C2 & C3 & C4 & C5 & \textbf{Total} \\ \hline
\textbf{Points} & 20 & 20 & 16 & 12 & 12 & 10 & 10 & \textbf{100} \\ \hline
\textbf{Score} & & & & & & & & \\[0.4em] \hline
\end{tabular}
\end{center}

\textbf{Rules.} This paper covers Weeks~1--7. [Closed book; no notes, no phones --- to be confirmed.] [A basic, non-programmable calculator is allowed --- to be confirmed.] Write your answers in the boxes and show every step of a calculation (formula $\to$ numbers $\to$ result with unit). \textbf{Course conventions:} context-switch time $0$; ties broken by arrival time, then by name; a smaller priority number means a higher priority; in Round Robin a process arriving at the moment a quantum expires enters the queue before the preempted process; TLB: $\text{EAT} = t_{\text{TLB}} + (2-\alpha)\,t_m$; holes are listed in address order and a request is carved from the beginning of the chosen hole; page frames start empty.

% ---------------------------------------------------------------------
\section*{Part A --- Multiple choice \hfill (10 $\times$ 2 = 20 points)}
Circle \emph{one} answer.
\begin{enumerate}[leftmargin=2.0em,label=A\arabic*.,itemsep=0.15em]
  \item The first successful time-sharing system, demonstrated at MIT in 1961, was \quad (a) Multics \quad (b) CTSS \quad (c) OS/360 \quad (d) MS-DOS
  \item In which OS structure do all services run as one program in kernel mode, so that it is fast but a buggy driver can crash the whole system? \quad (a) microkernel \quad (b) layered \quad (c) monolithic \quad (d) none of these
  \item In the \emph{child} process, \cmd{fork()} returns \quad (a) the child's PID \quad (b) $0$ \quad (c) the parent's PID \quad (d) $-1$
  \item Which transition does \emph{not} exist in the five-state process model? \quad (a) ready $\to$ running \quad (b) running $\to$ waiting \quad (c) waiting $\to$ running \quad (d) running $\to$ ready
  \item Which quantity does the CPU-scheduling algorithm influence most directly? \quad (a) the CPU burst length \quad (b) the I/O time \quad (c) the waiting time \quad (d) the size of the program
  \item Linux and Windows map user threads to kernel threads with the model \quad (a) many-to-one \quad (b) one-to-one \quad (c) many-to-many \quad (d) one-to-many
  \item If the Round Robin quantum is larger than every CPU burst, Round Robin behaves like \quad (a) SJF \quad (b) FCFS \quad (c) priority \quad (d) SRTF
  \item To use execution-time address binding, the system needs \quad (a) a linker \quad (b) a relocating loader \quad (c) an MMU \quad (d) compaction
  \item The TLB is \quad (a) a small, fast cache of page-table entries \quad (b) the page table itself \quad (c) a list of free frames \quad (d) a disk area for swapped pages
  \item A system is thrashing. The correct reaction of the OS is to \quad (a) admit more processes \quad (b) increase the time quantum \quad (c) replace LRU by OPT \quad (d) reduce the number of processes in memory
\end{enumerate}

% ---------------------------------------------------------------------
\section*{Part B --- Short answer \hfill (4 $\times$ 5 = 20 points)}
\begin{enumerate}[leftmargin=2.0em,label=B\arabic*.,itemsep=0.3em]
  \item Describe the two classic views of an operating system (\emph{extended machine} and \emph{resource manager}). Give one example of time multiplexing and one of space multiplexing. \pts{5}
  \ans{2.7cm}
  \item Compare user-level and kernel-level threads: give one advantage and one disadvantage of each. What happens in the many-to-one model when one thread makes a blocking system call? \pts{5}
  \ans{2.7cm}
  \item Two threads use semaphores $S$ and $Q$, both initialised to $1$: the threads execute\\[0.2em]
  \hspace*{1em}P0: \texttt{wait(S); wait(Q); ... signal(S); signal(Q);}\\
  \hspace*{1em}P1: \texttt{wait(Q); wait(S); ... signal(Q); signal(S);}\\[0.2em]
  Give an order of execution that leads to a deadlock, explain why, and state a simple rule that prevents it. \pts{5}
  \ans{2.7cm}
  \item Explain the difference between a hard link and a symbolic link. A user runs \cmd{ln a.txt h.txt}, \cmd{ln -s a.txt s.txt} and then \cmd{rm a.txt}. Can the data still be read through \texttt{h.txt}? Through \texttt{s.txt}? Why? \pts{5}
  \ans{2.7cm}
\end{enumerate}

\newpage
% ---------------------------------------------------------------------
\section*{Part C --- Problems \hfill (60 points)}
\begin{enumerate}[leftmargin=2.0em,label=C\arabic*.,itemsep=0.5em]
  \item \pts{16} Four processes arrive at a single CPU (times in ms):
  \begin{center}
  \begin{tabular}{@{}cccc@{}}
  \toprule
  \textbf{Process} & \textbf{Arrival} & \textbf{Burst} & \textbf{Priority} \\
  \midrule
  A & 0 & 6 & 3 \\ B & 2 & 5 & 2 \\ C & 3 & 2 & 1 \\ D & 5 & 4 & 4 \\
  \bottomrule
  \end{tabular}
  \end{center}
  For (a) non-preemptive SJF, (b) preemptive priority scheduling and (c) Round Robin with $q = 3$~ms, draw the Gantt chart and fill in the table. (d) Which of the three gives the smallest average waiting time?

  \ganttgrid{17}{(a)}
  \proctable
  \ganttgrid{17}{(b)}
  \proctable
  \ganttgrid{17}{(c)}
  \proctable
  \par\medskip (d) \hrulefill

  \newpage
  \item \pts{12}
  (a) A process has base register \texttt{0x6000} and limit register \texttt{0x0A00}. Translate the logical addresses \texttt{0x09FF}, \texttt{0x0A00} and \texttt{0x0150}. Give the legal range of physical addresses. \pts{4}
  \ans{2.6cm}
  (b) A segment table contains: segment $0$: base $1400$, limit $600$; segment $1$: base $6300$, limit $400$; segment $2$: base $4300$, limit $1100$; segment $3$: base $3200$, limit $1000$. Translate $\langle 2, 1099\rangle$, $\langle 1, 400\rangle$ and $\langle 3, 512\rangle$. \pts{4}
  \ans{2.4cm}
  (c) A paged system has $t_{\text{TLB}} = 20$~ns and $t_m = 100$~ns. Compute the effective access time for a hit ratio of $0.8$ and without a TLB. Which hit ratio gives an effective access time of $128$~ns? \pts{4}
  \ans{2.6cm}

  \item \pts{12} For the reference string $1, 0, 4, 2, 1, 0, 1, 2, 3, 1, 2, 3$ and $3$ initially empty frames, fill in the frame tables and count the page faults for FIFO, LRU and OPT.
  \par\medskip
  \frametable{FIFO \quad faults: \underline{\hspace{1.5cm}}}\par\medskip
  \frametable{LRU \quad faults: \underline{\hspace{1.5cm}}}\par\medskip
  \frametable{OPT \quad faults: \underline{\hspace{1.5cm}}}

  \newpage
  \item \pts{10} Memory has five holes (size@start, in KB, in address order): $650@200$, $450@1100$, $200@1800$, $550@2200$, $250@3000$. Processes A $= 270$, B $= 290$, C $= 330$ and D $= 510$~KB arrive in this order. For first-fit, best-fit and worst-fit, give the hole used by each process and the size of the hole that is left. Which requests cannot be satisfied, and why?
  \begin{center}
  \renewcommand{\arraystretch}{2.0}
  \begin{tabular}{|l|p{4.0cm}|p{4.0cm}|p{4.0cm}|}
  \hline
  \textbf{Request} & \textbf{First-fit} & \textbf{Best-fit} & \textbf{Worst-fit} \\ \hline
  A $= 270$ &&& \\ \hline
  B $= 290$ &&& \\ \hline
  C $= 330$ &&& \\ \hline
  D $= 510$ &&& \\ \hline
  \end{tabular}
  \end{center}
  \ans{2.0cm}

  \item \pts{10}
  (a) A file consists of $300$-byte records numbered from $0$; the file system uses $1024$-byte blocks. In which relative block, and at which offset inside it, does record $45$ start? Does record $58$ fit into a single block? \pts{4}
  \ans{3.0cm}
  (b) How many times is each letter printed by the following code (output to a terminal), and how many processes exist at the end? \pts{3}\\
  \texttt{fork(); printf("X\textbackslash n"); fork(); fork(); printf("Y\textbackslash n");}
  \ans{2.2cm}
  (c) A context switch takes $s = 20\ \mu$s and the time quantum is $q = 5$~ms. What percentage of the CPU time is lost to switching if every process uses its full quantum? \pts{3}
  \ans{2.0cm}
\end{enumerate}

\vfill
\begin{center}\small\color{muted}--- End of the sample paper. Good luck! ---\end{center}

\end{document}
